%% ma: L12-B15-0005
%% lop: 12
%% chuong: 5
%% bai: 15
%% dang: 12.15.1
%% loai: TN4
%% muc: NB
%% dap_an: "D"
%% nguon: "(NBV)-ĐỀ SỐ 6-MỖI NGÀY 1 ĐỀ THI 2025.pdf (D06, MỖI NGÀY 1 ĐỀ - Đề số 6), Phần I Câu 12"
%% kien_thuc: "Bài 15 (phương trình tham số của đường thẳng qua hai điểm), Bài 8 (tọa độ vectơ MN)"
%% trang_thai: chinh-thuc
%% ghi_chu: "Giáo viên đã duyệt 10/10/2026. Đáp án gốc D đúng. Có thể gộp với 12.15.1 khi duyệt"
\begin{ex}
Trong không gian $Oxyz$, cho hai điểm $M(4;-2;1)$ và $N(5;2;3)$. Đường thẳng $MN$ có phương trình là
\choice{$\begin{cases}x=-5+t\\y=2+4t\\z=3+2t\end{cases}$}{$\begin{cases}x=4-t\\y=-2-4t\\z=1+2t\end{cases}$}{$\begin{cases}x=4+t\\y=-2-4t\\z=1+2t\end{cases}$}{\True $\begin{cases}x=5-t\\y=2-4t\\z=3-2t\end{cases}$}
\loigiai{$\overrightarrow{MN}=(1;4;2)$ là một vectơ chỉ phương của $MN$; $(-1;-4;-2)$ cũng là một vectơ chỉ phương. Đường thẳng đi qua $N(5;2;3)$ nhận $(-1;-4;-2)$ làm vectơ chỉ phương có phương trình $x=5-t$, $y=2-4t$, $z=3-2t$. Các phương án A, B, C có vectơ chỉ phương hoặc điểm đi qua không phù hợp.}
\end{ex}
